\centerline{\bf \Huge Сравнения по модулю II}
\centerline{\bf \Huge Весёлая ферма}

\begin{flushright}
        \scriptsize{Хоть нам и грустно, но мы всё равно должны жить дальше,\\даже с этой грустью.\\ Канеки Кен}
\end{flushright}

\textbf{Ключевой факт.} Если $p$ - простое число и $a$ - целое число, не делящееся на $p$, то $a^{p-1}-1$ делится на $p$.

\textit{Примечание.} Все задачи из листочка кроме 8-2 нужно решать с помощью малой теоремы Ферма.

\problem{1}
Найдите остаток от деления

а) $2025^{2026}$ при делении на $2027$

б) $8^{900}$ на 29 

\problem{2}
Докажите, что $a^{p^2} \equiv a(\bmod p)$, где $p$ - простое число.

\problem{3}
Докажите, что если $p$ - простое число и $p>2$, то $7^p-5^p-2$ делится на $6 p$.

\problem{4}
Известно, что $a^{12}+b^{12}+c^{12}+d^{12}+e^{12}+f^{12}$ делится на 13 . Докажите, что $a b c d e f$ делится на $13^6$.

\problem{5}
\textbf{Первое число Кармайкла.}
Докажите, что $a^{561}-a$ делится на 561 для любого целого $a$.

\problem{6}
Докажите, что ни для какого целого $k$ число $k^2+k+1$ не кратно 101.

\problem{7}
Докажите, что для любого целого числа $a$

a) $a^5-a$ делится на 30 ;

б) $a^{17}-a$ делится на 510 ;

в) $a^{11}-a$ делится на 66.

\problem{8}
1) Докажите, что если $p-$ простое число и $1 \leq k \leq p-1$, то $C_p^k$ делится на $p$.

\noindent2) \textbf{Формула двоечника.} Пусть $p$ - простое число. Докажите, что $$(a+b)^p \equiv a^p+b^p(mod \ p)$$ для любых целых $a$ и $b$.